\scanpage{165}\mspage{23}
% archive sheet number: 165
Notons que tout quotient dans $\unsure{\mathscr X|\mathcal X}^{\wedge}$ d'un objet irréductible est irréductible, donc $w(x),w(a)$ sont des objets irréductibles. Cela montre que pour \inserted{$x$ dans $C$ et} $x'$ \deleted{dans} $S_{/W}$, \deleted{\unsure{on}} on a $x'$ dans $C$ i.e. $w(a)\cap w(x')\ne\varnothing$ ssi $w(x)\cap w(x')\ne\varnothing$.

[Lemme \inserted{trivial} : dans un objet \deleted{\unsure{irréduct.}} $W$ d'un topos $\mathscr X$, \deleted{\unsure{si}} on a \deleted{\unsure{deux}}\unsure{deux|trois} \unsure{sous-objets|objets} $W_1,W_2,W_3$ irréductibles, si l'un des deux rencontre les deux autres, ceux-ci se rencontrent.]

Cela montre que $C$ est un crible, puisque $x'\longrightarrow x$ implique $w(x')\subset w(x)$ donc $w(x')$ rencontre $w(a)$, donc aussi $w(\unsure{x|a})$. Mais \deleted{cela montre \inserted{alors} que $C'=S_{/W}\smallsetminus C$} le co-crible $C'=S_{/W}\smallsetminus C$ est également un crible, car si $x$ est dans $C'$ i.e. $w(x)$ ne rencontre pas $w(a)$, a fortiori pour tt $x'$ \deleted{\unsure{au}} dessous de $x$, pour lequel on aura $w(x')\subset w(x)$.

Ainsi, $S_{/W}$ est décomposé en \inserted{somme de} deux ouverts $C,C'$, non vides puisque l'un contient $a$, l'autre $b$. Par suite, $W$ est décomposé en somme de deux \inserted{$\ne$} sous-objets non vides, donc $W$ est non $0$-connexe.

On a prouvé 2.\unsure{8|9}, donc aussi \inserted{(revenant à 2.\unsure{6})} les implications vers $\mathrm{PS}1$ dans le diagramme page 20. Il reste à prouver

\[
\left\{\begin{aligned}
b&\Longrightarrow\mathrm{PS}2,\quad c'\Longrightarrow\mathrm{PS}2\\
\mathrm{PS}1+\mathrm{PS}2\ (\text{i.e. }\unsure{d|a})&\Longrightarrow\unsure{a'|b'}
\end{aligned}\right.
\]

\margin{NB finalement on ne \unsure{prouvera|prouve} pas $c'\Rightarrow\mathrm{PS}2$}
